\documentclass[12pt, twoside]{article}
\input{../preamble}

\fancyhead[LE]{\thepage}
\fancyhead[RO]{Markscheme}
\fancyhead[LO]{La Scuola d'Italia / Dr.\ Huson / IB Math 11 \\* 9 September 2026}

\begin{document}
\subsubsection*{1.2 Classwork: Arithmetic Sequences --- Markscheme}

\begin{enumerate}[itemsep=0.3cm]

%--- Problem 1: Recognising an arithmetic sequence, u_n and u_50 [8 marks] ---
\item[1.] [Maximum mark: 8]

The sequence $7,\; 11,\; 15,\; 19,\; \ldots$

\medskip

\textbf{(a)} $u_{1} = 7$ \hfill \textbf{A1}

\medskip

\textbf{(b)} The sequence is arithmetic. \hfill \textbf{A1}

Reason: the difference between consecutive terms is constant:
$11 - 7 = 15 - 11 = 19 - 15 = 4$. \hfill \textbf{R1}

\emph{Accept ``yes'', ``arithmetic'', or equivalent for the A1.}

\emph{R1 is awarded for any statement that the difference between
consecutive terms is the same (in words, or by showing at least two
equal differences). ``It goes up by 4 every time'' is sufficient.
Award the R1 if the reasoning is shown in a later part: a candidate who
states only that the sequence is arithmetic here, and then computes the
equal differences under (c), has given the reason and earns the R1.}

\medskip

\textbf{(c)} $d = 4$ \hfill \textbf{A1}

\medskip

\textbf{(d)} $u_{5} = 23$, $u_{6} = 27$ \hfill \textbf{A1}

\emph{Both values required for the single mark. Accept
``$23, 27$'' written as a continuation of the list.}

\emph{FT: award A1 for the next two terms obtained by adding the
candidate's own $d$ from (c) to $19$.}

\medskip

\textbf{(e)} $u_{n} = u_{1} + (n-1)d = 7 + (n-1)(4)$ \hfill \textbf{M1}

$u_{n} = 4n + 3$ \hfill \textbf{A1}

\emph{Accept the unsimplified form $7 + 4(n-1)$ for full marks.}

\emph{FT: award M1A1 for a correct expression formed from the
candidate's own $u_{1}$ and $d$.}

\medskip

\textbf{(f)} $u_{50} = 4(50) + 3 = 203$ \hfill \textbf{A1}

\emph{FT: award A1 for the value obtained by substituting $n = 50$
into the candidate's own expression from (e).}

\newpage

%--- Problem 2: Negative common difference; solve u_n = -50 [6 marks] ---
\item[2.] [Maximum mark: 6]

$u_{1} = 40$, $d = -3$.

\medskip

\textbf{(a)} $40,\; 37,\; 34,\; 31$ \hfill \textbf{A1}

\emph{All four terms required for the mark.}

\medskip

\textbf{(b)} $u_{20} = 40 + (20 - 1)(-3) = 40 - 57$ \hfill \textbf{M1}

$u_{20} = -17$ \hfill \textbf{A1}

\emph{M1 is for any correct use of $u_{n} = u_{1} + (n-1)d$ with
$n = 20$, however the arithmetic is then arranged; a GDC sequence
table showing the $20^{\text{th}}$ term also earns the M1.}

\medskip

\textbf{(c)} $40 + (n - 1)(-3) = -50$ \hfill \textbf{M1}

$(n - 1)(-3) = -90$

$n - 1 = 30$ \hfill \textbf{A1}

$n = 31$ \hfill \textbf{A1}

\emph{Accept the equivalent set-up $43 - 3n = -50$ leading to
$3n = 93$.}

\emph{The A1 for $n - 1 = 30$ is for any correct intermediate line
that isolates $n$ (e.g.\ $-3n = -93$); it does not have to be written
in this exact form.}

\emph{A candidate who states $n = 31$ from a GDC table with no
algebraic working earns the final A1 only, since the question asks
for algebraic working.}

\end{enumerate}
\end{document}